\section{Results and Evaluation}\label{sec:results}

We evaluate HospitAI on the synthetic cohort of \NPATIENTS{} patients. The goal
is not to claim clinical validity on synthetic data, but to show that the
pipeline produces \emph{interpretable, internally consistent} structure and that
each stage feeds the next as intended.

\subsection{Clustering quality}
With $k=\KCLUSTERS$, K-Prototypes partitions the cohort into clinically
distinguishable groups. Table~\ref{tab:metrics} reports internal validity indices
over both the numerical subspace and the mixed (Gower) space. The per-feature
$\eta^2$ and Cram\'er's $V$ scores identify which attributes drive the
separation --- in our runs the dominant care-plan category, condition burden and
the severity score are among the strongest discriminators --- which gives the
clustering an interpretable rationale rather than an opaque label.

\begin{table}[!t]
  \caption{Internal cluster validity ($k=\KCLUSTERS$).}
  \label{tab:metrics}
  \centering
  \begin{tabular}{lc}
    \toprule
    \textbf{Metric} & \textbf{Value} \\
    \midrule
    Silhouette (numerical)        & \TODO{fill from metrics} \\
    Silhouette (Gower, mixed)     & \TODO{fill} \\
    Davies--Bouldin (lower better)& \TODO{fill} \\
    Calinski--Harabasz (higher better) & \TODO{fill} \\
    \bottomrule
  \end{tabular}
\end{table}

\begin{figure}[!t]
  \centering
  % Export the FAMD scatter persisted in cluster_meta.pca_image.
  \fbox{\parbox{0.9\columnwidth}{\centering\vspace{2.0cm}
  \textbf{FAMD 2-D cluster scatter (placeholder)}\vspace{2.0cm}}}
  \caption{Two-dimensional FAMD projection of the mixed feature space, coloured
  by cluster. Replace with the figure stored during the clustering run.}
  \label{fig:famd}
\end{figure}

\subsection{Cluster-specific trajectories}
The central qualitative result is the contrast between global and per-cluster
sequence mining. Table~\ref{tab:patterns} illustrates the intended finding: a
trajectory can be dominant inside one cluster while remaining infrequent across
the whole population, exposing a care progression that is specific to that
patient group and would be diluted by global mining alone. The exact patterns and
support values should be filled in from the persisted pattern tables.

\begin{table}[!t]
  \caption{Example trajectories: per-cluster vs.\ global support.
  Fill from the mined pattern tables.}
  \label{tab:patterns}
  \centering
  \footnotesize
  \begin{tabular}{p{0.50\columnwidth}cc}
    \toprule
    \textbf{Pattern} & \textbf{Cluster} & \textbf{Global} \\
    \midrule
    \texttt{C:diabetes $\rightarrow$ M:metformin $\rightarrow$ P:foot\_exam}
      & \TODO{\%} & \TODO{\%} \\
    \TODO{pattern 2} & \TODO{\%} & \TODO{\%} \\
    \TODO{pattern 3} & \TODO{\%} & \TODO{\%} \\
    \bottomrule
  \end{tabular}
\end{table}

\subsection{Interpretation layer}
Qualitatively, the LLM layer renders each cluster as a titled, plain-language
profile with a risk level and a recommended action, and each trajectory as a
short clinical narrative. The text-to-SQL agent answers natural-language
questions (e.g., \emph{``list patients with more than three emergency encounters
in the past year''}) by emitting a validated read-only query whose text is shown
alongside the answer. Because validation rejects non-\texttt{SELECT} statements
and unknown schema elements before execution, malformed or unsafe generations are
caught rather than run. \TODO{optionally add a small example screenshot or a
table of representative questions, generated SQL, and whether validation passed
on the first attempt.}
